\section{Arithmetic setup and computational method}\label{sec:setup}
\subsection{Primes, residue symbols and normalisation}
Put \(\omega=e^{2\pi i/3}\) and \(\zeta_6=1+\omega\).
We use the repository's primary convention
\(\rho=a+b\omega\), \(a\equiv2\pmod3\), \(b\equiv0\pmod3\),
with rational prime norm \(p=\Nm(\rho)\le X\).
There are two such chosen degree-one primes above each rational
\(p\equiv1\pmod3\). Define \(D=|\PP(X)|\).
For a denominator prime \(\sigma\), the symbol
\(\chi_\sigma(u)=(u/\sigma)_6\) is the sextic residue symbol,
and is zero whenever \(\sigma\mid u\).
The residue matrix has \emph{numerator row} and \emph{denominator
column}:
\[
 (\rX)_{\rho,\sigma}=\chi_\sigma(\rho),\qquad
 (\rX)_{\rho,\rho}=0.
\]

Write \(e(z)=e^{2\pi iz}\), and for \(p=\Nm(\sigma)\) define
the rational-normalized Gauss sum
\[
 \gtilde(\sigma)=p^{-1/2}\sum_{a=1}^{p-1}
       \chi_\sigma(a)e(a/p).
\]
The rational additive character is essential; it differs from the
trace character used in some composite-modulus comparisons.

A twist class \(\kappa=\prod_\pi\pi^{e_\pi}\) has finitely
many distinct factors, \(1\le e_\pi\le5\), taken modulo sixth powers.
The factors may be split or inert, subject to the original
coprimality convention. Put
\[
 n=\Nm_6(\kappa)=\prod_\pi\Nm(\pi)^{e_\pi},\qquad
 w_\kappa(\sigma)=\overline{\gtilde(\sigma)}\chi_\sigma(\kappa).
\]
Define, where denominators are nonzero,
\begin{align}
 r_\kappa(X)&=\frac{\|\rX w_\kappa\|}
                  {\sqrt D\,\|w_\kappa\|},&
 \epsilon_\kappa(X)&=\frac{r_\kappa(X)^2-1}{r_1(X)^2-1},
 \label{eq:r}\\
 s_\kappa(X)&=n^{1/6}\epsilon_\kappa(X),&
 t&=\frac nX.\label{eq:st}
\end{align}
A value \(r=1\) is a natural random-phase benchmark, but a finite
comparison must use the symmetry-matched random control, not an
unrestricted independent random matrix.

\fig{primary_lattice}{.80}{Lattice plot of sample primary Eisenstein
prime representatives in the RS-S019 class list, from
\texttt{results/rs\_s019\_moduli\_2026-10-26.json}
(\texttt{30965467}). This is illustrative of indexing, not
a density measurement.}{lattice}

\subsection{Experiment design, controls and computational validation}
The programme used committed pre-data predictions with explicit
tolerances. Discovery, new-class holdouts and new-scale holdouts
are distinct. Failed tests were recorded without retrospectively
changing their criteria. In particular, the RS-S015
norm-only proposal and RS-S017 per-class fifth-power test failed;
neither is described as a success merely because a weaker
grouped statistic was subsequently supported.

The null has the same matrix dimension and reciprocity/conjugation
structure as the observed ensemble. Native finite-field and
matrix-free implementations were cross-checked against existing
engine outputs; the RS-S021 reproducibility file records a
relative difference of at most \(5.27\times10^{-15}\)
at its specified validation scale.
% Source: reciprocity_spectra/results/rs_s021_large_x_2026-10-26.json, validation.X320000.max_relative_r_difference.
This is an agreement of two numerical implementations, not a
theorem about an infinite family.

At large \(X\), a matrix-free native engine evaluates matrix--vector
products without materializing the full dense \(D\times D\) array.
Its code, run parameters, platform and control seeds are preserved
with the result JSON. We do not treat repeated views of the same
data as independent replications.

\fig{prereg_timeline}{.96}{Chronology of finite prediction outcomes:
RS-S015 failure; RS-S016 exploratory regime study; RS-S017 mixed
pass/fail; RS-S019, S020 and S021 registered passes of their specified
decisions. Source score blocks in the six corresponding
\texttt{results/rs\_s0**.json} files at \texttt{30965467}.
Timeline is schematic, not a new experiment.}{timeline}

\begin{table}[htbp]\centering\small
\begin{tabular}{@{}llll@{}}\toprule
Session & Frozen test & Outcome & Preserved caveat\\\midrule
S015 & PROV-RS14 P1--P3 & FAIL/PASS/FAIL & precursor refuted\\
S017 & Q1,Q2,Q4 & PASS & fresh classes\\
S017 & Q3 & FAIL & per-class \(r<1\) clause\\
S019 & prime/all-moduli & PASS & finite comparison\\
S020 & H-minus & PASS & grouped fifth classes\\
S021 & new-scale A, B & PASS & finite \(X\le4\cdot10^6\)\\
\bottomrule\end{tabular}
\caption{Selected frozen prediction outcomes. The complete
criteria and provenance reside in
\texttt{research/CONJECTURE\_REGISTER.md} and
\texttt{reciprocity\_spectra/results/} (\texttt{30965467}).}
\label{tab:pred}
\end{table}
